\documentclass[12pt,oneside,a4paper]{report}

% ----------------------- USTAWIENIA JĘZYKA I CZCIONEK -----------------------
\usepackage[polish]{babel}
\usepackage[T1]{fontenc}
\usepackage[utf8]{inputenc}
\usepackage{lmodern}
\usepackage{microtype}

% ----------------------- UKŁAD STRONY --------------------------------------
\usepackage[a4paper,margin=2.5cm,includefoot]{geometry}
\usepackage{setspace}
\onehalfspacing
\usepackage{indentfirst}

% ----------------------- HIPERŁĄCZA ----------------------------------------
\usepackage[hidelinks]{hyperref}
\usepackage{bookmark}

% ----------------------- CYTOWANIA I BIBLIOGRAFIA ---------------------------
\usepackage{csquotes}
\usepackage[backend=biber,style=authoryear,sorting=nyt]{biblatex}
\addbibresource{references.bib}

% ----------------------- DANE DOKUMENTU ------------------------------------
\newcommand{\autor}{Jan Zakroczymski}
\newcommand{\tytul}{FAscnn++: Rozwinięcie architektury Fast-SCNN o mechanizm Fast Attention i efektywny moduł fuzji cech}
\newcommand{\uczelnia}{UKSW - Uniwersytet Kardynała Stefana Wyszyńskiego}
\newcommand{\rok}{2025}

\begin{document}

\begin{titlepage}
  \begin{center}
    {\Large \uczelnia \\[0.5cm]}
    {\Large Sprawozdanie z postępów pracy (rozszerzony abstrakt)}\\[1.5cm]
    {\LARGE \textbf{\tytul}}\\[1.5cm]
    {\large Autor: \autor}\\[0.2cm]
    {\large Rok: \rok}\\[2cm]
  \end{center}
\end{titlepage}

\tableofcontents

\chapter{Wstęp}
Semantyczna segmentacja obrazu polega na przypisaniu każdemu pikselowi etykiety
klasy opisującej znaczenie fragmentu sceny. W odróżnieniu od klasyfikacji i
detekcji obiektów, segmentacja semantyczna wymaga jednoczesnego zachowania
szczegółów geometrycznych oraz spójność kontekstową w skali całego obrazu.
W praktyce oznacza to, że algorytm musi rozumieć zarówno granice obiektów
(np. krawędzie pojazdu), jak i relacje pomiędzy elementami sceny (np. fakt, że
chodnik graniczy z jezdnią). Jest to zadanie obliczeniowo wymagające, ale
niezwykle istotne w systemach wspomagania kierowcy, autonomicznych robotach
mobilnych, diagnostyce medycznej czy monitoringu wizyjnym.

Rozwój sieci konwolucyjnych oraz technik uczenia głębokiego spowodował szybki
postęp w jakości segmentacji. Modele takie jak FCN, U-Net, SegNet, DeepLab czy
PSPNet ustanowiły standardy dokładności na popularnych zbiorach danych
\parencite{long2015fully,ronneberger2015u,badrinarayanan2017segnet,chen2018deeplab,zhao2017pspnet}.
Jednocześnie wzrosła złożoność modeli, co utrudnia ich zastosowanie w systemach
wbudowanych i urządzeniach o ograniczonych zasobach. Dla aplikacji w czasie
rzeczywistym krytyczne są trzy czynniki: liczba operacji obliczeniowych,
zapotrzebowanie na pamięć oraz opóźnienie inferencji.

W odpowiedzi na te potrzeby rozwijane są architektury lekkie, ukierunkowane na
szybkie przetwarzanie. ENet, Fast-SCNN czy BiSeNet wprowadzają uproszczone bloki
konwolucyjne oraz modułowe podejście do ekstrakcji cech, osiągając wysoką
prędkość na GPU i CPU \parencite{paszke2016enet,mehta2019fast,yu2018bisenet}.
Równolegle rozwijają się modele wykorzystujące uwagę i transformatory, które
poprawiają modelowanie kontekstu sceny \parencite{vaswani2017attention,dosovitskiy2020vit,xie2021segformer}.
Wyzwaniem pozostaje połączenie szybkiego przetwarzania z silnym modelowaniem
zależności globalnych.

Niniejsze sprawozdanie stanowi rozszerzony abstrakt pracy obejmujący jej pierwsze
dwa rozdziały. Przedstawia kontekst problemu, motywację, tezę i metodologię
badawczą, a także przegląd literatury w formie bibliografii. Dokument opisuje
postępy nad autorską architekturą FAscnn++ (Fast Atencion SCNN) rozwijaną w
ramach projektu FAscnn++ Fast Segmentation.

\section{Zakres i cele sprawozdania}
Celem sprawozdania jest uporządkowanie dotychczasowych prac i wskazanie jasnych
kierunków dalszego rozwoju. Dokument podsumowuje fundamenty teoretyczne oraz
uzasadnia potrzebę projektowania szybkiej architektury segmentacyjnej na CPU.
W dalszej części zostanie przedstawiona teza oraz metodologia, które prowadzą
od sformułowania problemu do weryfikacji hipotezy badawczej.

\section{Ogólny opis projektu}
Projekt FAscnn++ Fast Segmentation to implementacja kilku modeli segmentacji
semantycznej, w tym ENet, jego modyfikacji oraz autorskiej architektury FAscnn++.
Kod zorganizowano w postaci modułu treningowego i testowego w języku Python
(PyTorch), z zestawem skryptów do treningu i ewaluacji. Jako podstawowy zbiór
danych przyjęto Cityscapes \parencite{cordts2016cityscapes}, który jest
standardem w badaniach segmentacji ulicznej.

Z perspektywy użyteczności, architektura FAscnn++ ma docelowo umożliwiać działanie na
urządzeniach o ograniczonej mocy obliczeniowej przy zachowaniu wysokiej jakości
segmentacji. Projekt zakłada modułowy charakter sieci, umożliwiający łatwe
porównanie z klasycznymi modelami oraz wprowadzanie wariantów architektury.

\section{Cele naukowe i inżynierskie}
Cele pracy mają dwojaki charakter. Z perspektywy naukowej celem jest zbadanie,
czy uproszczony mechanizm uwagi może poprawić jakość segmentacji w lekkich
architekturach. Oznacza to nie tylko zaprojektowanie nowego modułu, ale także
analizę jego wpływu na mIoU i stabilność predykcji w różnych warunkach.

Z perspektywy inżynierskiej celem jest stworzenie praktycznego modelu, który
można uruchomić na sprzęcie bez dedykowanego GPU. W konsekwencji rozwiązanie
musi być łatwe do trenowania, reprodukowalne oraz kompatybilne z typowymi
pipeline'ami przetwarzania obrazów. Te cele determinują zarówno wybór architektury,
jak i zakres eksperymentów.

\section{Znaczenie problemu w praktyce}
Zadania segmentacji semantycznej pojawiają się wszędzie tam, gdzie system musi
zrozumieć strukturę sceny w sposób ciągły i przestrzenny. W ruchu drogowym
segmentacja umożliwia rozróżnianie jezdni, chodników, znaków i uczestników ruchu.
W robotyce pozwala planować trajektorie ruchu na podstawie mapy przeszkód i
regionów przejezdnych. W diagnostyce medycznej segmentacja odpowiada za
precyzyjne wydzielanie tkanek, co wspiera planowanie terapii i analizę
patologii. W monitoringu wizyjnym kluczowe jest szybkie wykrywanie obszarów
zainteresowania i obiektów, które pojawiają się na obrazie.

We wszystkich tych zastosowaniach systemy działają w warunkach silnych
ograniczeń czasowych. Długie opóźnienie przetwarzania oznacza ryzyko błędnej
reakcji w krytycznym momencie. Dlatego w praktyce liczy się nie tylko sama
jakość segmentacji, lecz także stabilność działania, przewidywalność czasu
inferencji oraz łatwość wdrożenia na różnorodnych platformach sprzętowych.

\section{Wyzwania obliczeniowe}
Współczesne modele segmentacji korzystają z głębokich sieci, które generują
wiele map cech o wysokiej rozdzielczości. Kluczowym wyzwaniem jest utrzymanie
balansu pomiędzy szczegółowością wyników a kosztami obliczeniowymi. Wysoka
rozdzielczość ułatwia rekonstrukcję granic obiektów, ale prowadzi do dużej
liczby operacji i wymaganej pamięci. Z drugiej strony agresywny downsampling
zmniejsza koszt, lecz powoduje utratę drobnych detali.

Kolejnym problemem jest dostępność zasobów. W zastosowaniach wbudowanych
często brakuje GPU, a modele muszą działać na wielordzeniowych CPU lub
energooszczędnych akceleratorach. W takim środowisku kluczowe jest ograniczenie
złożoności obliczeniowej oraz unikanie operacji o wysokim koszcie pamięciowym,
takich jak pełna uwaga kwadratowa.

Nie można pominąć aspektu danych. Wysokiej jakości adnotacje pikselowe są
drogie w pozyskaniu, co ogranicza rozmiar zbiorów treningowych. W praktyce
model musi efektywnie wykorzystywać ograniczony zbiór danych, a procedury
augmentacji powinny wzmacniać uogólnianie bez wprowadzania sztucznych artefaktów.
Z perspektywy projektu oznacza to konieczność starannego zaprojektowania
strategii treningowej oraz kontrolowania ryzyka przeuczenia.

\section{Wymagania real-time}
System działający w czasie bliskim rzeczywistemu musi zapewniać stabilny czas
inferencji. Oprócz średniej liczby FPS istotna jest zmienność czasu
przetwarzania pomiędzy próbkami. Dla aplikacji sterujących pojazdami lub
robotami, priorytetem jest przewidywalność opóźnienia i niezawodność modelu.
Projekt FAscnn++ uwzględnia te wymagania, proponując architekturę, która minimalizuje
wahania czasu przetwarzania, a jednocześnie zachowuje zdolność do modelowania
kontekstu sceny.

\section{Definicje i zakres pojęciowy}
W obrębie przetwarzania obrazów wyróżnia się kilka pokrewnych zadań. Segmentacja
semantyczna dotyczy przypisania etykiety klasowej każdemu pikselowi i nie
rozróżnia instancji tego samego typu obiektów. Segmentacja instancyjna
wydziela poszczególne obiekty w ramach jednej klasy, a segmentacja panoptyczna
łączy oba podejścia w jeden system \parencite{he2017maskrcnn,kirillov2019panoptic}.
W niniejszej pracy koncentruje się na segmentacji semantycznej, ponieważ stanowi
ona naturalny punkt wyjścia dla zastosowań real-time.

Ważne jest rozróżnienie pomiędzy metrykami jakości i wydajności. Metryki
jakości, takie jak mIoU, opisują zgodność maski predykcyjnej z prawdą
referencyjną. Metryki wydajności, takie jak FPS czy opóźnienie, określają
prędkość inferencji. W praktyce oba aspekty muszą być uwzględnione jednocześnie,
ponieważ wysoka jakość bez wymaganej prędkości nie spełnia założeń aplikacyjnych.

\section{Kontekst badawczy}
W literaturze można zauważyć rosnący trend rozwijania architektur lekkich oraz
technik optymalizacyjnych, które pozwalają wdrażać modele w systemach
wbudowanych. Coraz częściej publikacje raportują nie tylko mIoU, ale również
parametry wydajności, liczbę operacji FLOPs czy zapotrzebowanie na pamięć.
W tym kontekście projekt FAscnn++ wpisuje się w nurt badań nad segmentacją
zorientowaną na zastosowania praktyczne i rzeczywiste ograniczenia sprzętowe.

Istotne jest także przesunięcie w stronę modeli, które potrafią balansować
lokalny kontekst z globalnymi zależnościami. Tradycyjne modele konwolucyjne
opierają się na lokalnych filtrach, a zrozumienie zależności dalekiego zasięgu
wymaga głębokich sieci i wieloskalowych mechanizmów. Rozwiązania oparte o uwagę
rozszerzają te możliwości, ale często kosztem złożoności obliczeniowej.
W efekcie istnieje realna potrzeba rozwiązań hybrydowych, które łączą zalety
obu podejść.

\chapter{Motywacja}
Motywacja pracy wynika z łączenia dwóch sprzecznych wymagań: wysokiej jakości
segmentacji oraz niskiego opóźnienia przetwarzania. Klasyczne modele zapewniają
wysoką dokładność kosztem dużej liczby parametrów i kosztu obliczeniowego.
W zastosowaniach wbudowanych, takich jak autonomiczne drony, roboty magazynowe
czy systemy ADAS, model musi działać na CPU lub słabym akceleratorze i reagować
niemal natychmiast. Ograniczenia te prowadzą do poszukiwania rozwiązań, które
wykorzystują wiedzę z literatury, ale upraszczają architekturę i optymalizują
przetwarzanie.

Istotnym problemem w lekkich modelach jest utrata kontekstu globalnego. Modele
real-time często operują na wysoce zredukowanych mapach cech, co ogranicza
zdolność rozróżniania obiektów o podobnej teksturze lub barwie. Mechanizmy
uwagi są skuteczne w modelowaniu relacji długozasięgowych, ale klasyczne
implementacje uwagi są kosztowne obliczeniowo. Stwarza to potrzebę opracowania
uproszczonego, szybkiego mechanizmu uwagi, który można integrować z lekkimi
blokami konwolucyjnymi.

Dodatkowa motywacja wynika z rosnącego znaczenia segmentacji w zastosowaniach
praktycznych. Mapowanie drzew, detekcja infrastruktury drogowej, a także analiza
medyczna lub przemysłowa korzystają z segmentacji jako elementu kluczowego
pipeline'u. Dostęp do wydajnych modeli na CPU umożliwia wdrożenia w terenie,
gdzie nie ma dostępu do wyspecjalizowanego sprzętu.

W kontekście projektu FAscnn++ priorytetem jest stworzenie architektury, która łączy
zalety wielogałęziowego przetwarzania (multi-branch) z uwagą przyjazną dla CPU.
Rozwiązanie ma być jednocześnie proste w implementacji, łatwe do reprodukcji
i kompatybilne z istniejącymi pipeline'ami treningowymi.

\section{Problemy badawcze}
Motywacja prowadzi do następujących problemów badawczych:
\begin{itemize}
  \item Jak zminimalizować koszt obliczeniowy bez drastycznej utraty jakości?
  \item Czy uproszczony mechanizm uwagi może zwiększyć skuteczność lekkiej sieci?
  \item Jak zaprojektować modularną architekturę, która pozwala na szybkie testy
        wariantów i porównania z modelami bazowymi?
  \item Jak zbalansować liczbę kanałów, głębokość sieci oraz rozmiar wejścia tak,
        aby model pozostawał użyteczny na CPU?
\end{itemize}

\section{Oczekiwane korzyści}
Oczekiwanym efektem pracy jest opracowanie architektury, która pozwoli uzyskać
konkurencyjną jakość segmentacji na Cityscapes przy istotnie niższych kosztach
obliczeniowych niż standardowe modele. Dodatkową wartością jest stworzenie
pakietu kodu umożliwiającego reprodukcję wyników, trening oraz testy porównawcze.

\section{Założenia projektowe}
Projekt zakłada, że architektura FAscnn++ będzie kompatybilna z typowymi pipeline'ami
uczenia głębokiego. Oznacza to wykorzystanie standardowych warstw konwolucyjnych
i prostych modułów, które są dobrze wspierane przez biblioteki uczenia maszynowego.
Istotne jest także unikanie nadmiernej zależności od specyficznych funkcji
sprzętowych, co mogłoby utrudnić przenoszenie modelu między platformami.

Zakłada się również możliwość stosowania metod optymalizacji po treningu, takich
jak kwantyzacja lub przycinanie wag, które dodatkowo zmniejszają koszt
inferencji \parencite{jacob2018quantization,han2015deepcompression}.
Uwzględnienie tych narzędzi w projektowaniu architektury pozwala utrzymać
przewagę w realnych wdrożeniach.

\section{Luka badawcza}
Istnieje wyraźna luka pomiędzy modelami o bardzo wysokiej jakości a modelami
o wysokiej wydajności. Wiele architektur real-time skupia się na uproszczeniu
przepływu obliczeń, ale rezygnuje z mechanizmów, które w nowoczesnych modelach
zapewniają bogaty kontekst. Z drugiej strony, modele wykorzystujące uwagę
osiągają lepsze wyniki, lecz koszt obliczeniowy często wyklucza je z zastosowań
na CPU. FAscnn++ ma wypełnić tę lukę poprzez dostarczenie efektywnego mechanizmu
uwagi, który nie blokuje zastosowań w czasie rzeczywistym.

\section{Uwarunkowania wdrożeniowe}
Wdrożenie modelu segmentacji w systemach produkcyjnych wymaga uwzględnienia
ograniczeń energetycznych, kosztów sprzętowych oraz niezawodności. W wielu
przypadkach docelowe platformy to komputery jednopłytowe, kontrolery wbudowane
lub jednostki o niskim poborze mocy. Oznacza to, że nawet niewielka poprawa
wydajności może mieć duże znaczenie dla realnych zastosowań, np. wydłużenia
czasu pracy na baterii lub obniżenia kosztów infrastruktury.

Ponadto w realnych systemach ważna jest łatwość integracji. Architektury, które
nie wymagają specjalistycznych operatorów i pozwalają na uruchomienie na
standardowych bibliotekach, są znacznie łatwiejsze do wdrożenia. Projekt FAscnn++
ma na celu dostarczenie takiej architektury, która nie tylko spełnia wymagania
wydajnościowe, ale również jest przyjazna dla procesów inżynierskich i testów.

\chapter{Teza i metodologia}
\section{Teza pracy}
Teza pracy brzmi następująco: \emph{możliwe jest zaprojektowanie lekkiej
architektury segmentacji semantycznej, która dzięki wielogałęziowemu przetwarzaniu
oraz uproszczonemu mechanizmowi uwagi osiąga konkurencyjną jakość segmentacji
przy czasie inferencji umożliwiającym działanie w czasie bliskim rzeczywistemu
na CPU}.

Teza zakłada, że odpowiednio skonstruowany blok uwagi, o liniowej złożoności
obliczeniowej względem liczby pikseli, może poprawić modelowanie kontekstu bez
znacznego wzrostu kosztu. Dodatkowo zakłada, że architektura wielogałęziowa
pozwoli połączyć cechy lokalne i globalne, minimalizując utratę informacji.

\section{Hipotezy pomocnicze}
Weryfikacja tezy zostanie wsparta przez następujące hipotezy pomocnicze:
\begin{itemize}
  \item Zastosowanie głębokościowo-separowalnych konwolucji obniża koszt
        obliczeniowy przy zachowaniu reprezentacji cech \parencite{howard2017mobilenets,chollet2017xception}.
  \item Użycie modułów downsamplingu oraz fuzji cech w stylu Fast-SCNN zwiększa
        wydajność i pozwala na kontrolowaną utratę rozdzielczości
        \parencite{mehta2019fast}.
  \item Mechanizm uwagi o zredukowanej złożoności (np. liniowej) poprawia
        segmentację obiektów o złożonych kształtach bez znaczącego wzrostu czasu
        inferencji \parencite{katharopoulos2020transformers,wang2020linformer}.
\end{itemize}

\section{Metodologia badawcza}
Metodologia ma charakter projektowo-eksperymentalny i składa się z kilku etapów:

\subsection{Analiza literatury i wymagań}
Pierwszym krokiem jest przegląd literatury, obejmujący modele segmentacji oraz
techniki optymalizacji obliczeń. Na tej podstawie wyodrębniane są cechy, które
mają kluczowe znaczenie dla szybkości i jakości segmentacji, takie jak
downsampling, fuzja cech oraz uwaga. Równolegle analizowane są wymagania
sprzętowe, z naciskiem na inferencję na CPU.

\subsection{Projekt architektury FAscnn++}
Projekt architektury zakłada użycie kilku gałęzi przetwarzania. Jedna gałąź
ekstrahuje cechy lokalne na wyższej rozdzielczości, druga buduje reprezentacje
globalne na silnie zredukowanej siatce, a trzecia integruje informacje
kontekstowe za pomocą szybkiej uwagi. W praktyce architektura wykorzystuje
bloki typu bottleneck, głębokościowo-separowalne konwolucje oraz moduły fuzji
cech podobne do tych z Fast-SCNN. Mechanizm uwagi jest uproszczony tak, aby
złożoność obliczeniowa była zbliżona do liniowej.

Kluczowym elementem jest sposób połączenia gałęzi. W FAscnn++ zakłada się, że
informacje lokalne powinny zachować precyzję granic, natomiast informacje
globalne powinny stabilizować klasyfikacje w regionach o słabych teksturach.
Warstwa fuzji ma łączyć obie reprezentacje w sposób umożliwiający łatwą
adaptację do różnych rozdzielczości wejściowych. W praktyce oznacza to
zastosowanie operacji interpolacji oraz połączeń resztkowych.

\subsection{Implementacja i pipeline treningowy}
Model jest implementowany w PyTorch. Pipeline treningowy obejmuje:
\begin{itemize}
  \item przygotowanie danych Cityscapes z odpowiednim podziałem na zbiory,
  \item trenowanie modeli bazowych (ENet, Fast-SCNN) jako punktów odniesienia,
  \item trenowanie wariantów FAscnn++ w identycznych warunkach, aby zachować
        porównywalność wyników,
  \item logowanie metryk oraz zapisywanie najlepszych wag.
\end{itemize}
Ważnym elementem jest utrzymanie jednolitych hiperparametrów, takich jak
rozmiar batcha, liczba epok i strategie optymalizacji, co umożliwia sprawiedliwe
porównanie.

\subsection{Ewaluacja i metryki}
Ocena modeli odbywa się na podstawie standardowych metryk segmentacji:
\begin{itemize}
  \item \textbf{mIoU} (mean Intersection over Union) jako główna metryka jakości,
  \item \textbf{Pixel Accuracy} do oceny ogólnej poprawności klasyfikacji pikseli,
  \item \textbf{FPS oraz czas inferencji} mierzone na CPU i GPU,
  \item \textbf{Liczba parametrów i FLOPs} jako wskaźniki złożoności.
\end{itemize}
Dodatkowo przewiduje się analizę jakościową, tj. ocenę wizualną wyników na
przykładowych obrazach, ze szczególnym uwzględnieniem granic obiektów.

\subsection{Eksperymenty porównawcze i ablacjne}
Planowane są eksperymenty porównawcze z modelami referencyjnymi. Dodatkowo
przewiduje się eksperymenty ablacjne, które ocenia wpływ poszczególnych modułów
(np. uwagi, dodatkowych gałęzi) na wynik końcowy. Pozwoli to oszacować, które
elementy architektury przynoszą największą poprawę przy minimalnym koszcie.

\section{Kryteria sukcesu}
Za sukces uznaje się uzyskanie mIoU na Cityscapes na poziomie porównywalnym z
lekkimi modelami literaturowymi przy jednoczesnym osiągnięciu czasu inferencji
na CPU pozwalającym na przetwarzanie w czasie bliskim rzeczywistemu. Dodatkowym
kryterium jest stabilność trenowania oraz możliwość wdrożenia w praktycznym
pipeline'ie.

\section{Plan eksperymentów}
Plan eksperymentów obejmuje kilka etapów. W pierwszej kolejności trenowane są
modele bazowe z literatury, tak aby uzyskać punkty odniesienia dla metryk
jakości i wydajności. Następnie trenowane są warianty FAscnn++ z różnymi konfiguracjami
liczby gałęzi, szerokości kanałów oraz ustawień mechanizmu uwagi. W trzecim
etapie wykonywane są eksperymenty ablacjne, w których usuwa się wybrane elementy
architektury (np. moduł uwagi lub jedna z gałęzi), aby zbadać ich wpływ na wynik.

\section{Zasady porównań i raportowania}
Porównania będą prowadzone przy zachowaniu jednorodnych warunków treningu i
ewaluacji. W szczególności wszystkie modele będą trenowane przez tę samą liczbę
epok, na identycznym podziale danych, z tym samym schematem augmentacji.
Raportowanie wyników obejmie tabele metryk jakości oraz wykresy zależności
pomiędzy mIoU a kosztem obliczeniowym. Dodatkowo przewiduje się prezentację
przykładowych map segmentacji w celu zilustrowania różnic w jakości granic.

\section{Reprodukowalność i narzędzia}
Projekt opiera się na PyTorch oraz zestawie skryptów uruchamiających trening
i testy. Dane wejściowe są przygotowywane zgodnie z typowym podziałem Cityscapes,
a konfiguracje treningowe zapisywane są w logach. Dla zapewnienia powtarzalności
wyników przewiduje się jawne ustawienie ziarna losowości oraz raportowanie
kluczowych parametrów, takich jak rozmiar wejścia, liczba epok i typ optymalizatora.

\section{Przygotowanie danych i augmentacje}
Zbiór Cityscapes zawiera obrazy o wysokiej rozdzielczości, co w praktyce wymaga
skalowania lub wycinków podczas treningu. W projekcie planuje się stosowanie
augmentacji geometrycznych (losowe przycięcia, odbicia poziome) oraz
augmentacji fotometrycznych (zmiany jasności i kontrastu), aby poprawić
uogólnianie modelu. Szczególne znaczenie mają augmentacje, które nie naruszają
struktury obiektów o znaczeniu semantycznym, np. znaków drogowych.

W celu kontroli porównywalności wszystkie modele będą trenowane na tym samym
zestawie transformacji. Pozwoli to jednoznacznie przypisać ewentualne różnice
w wynikach do zmian architektury, a nie do różnych procedur przygotowania danych.

\section{Profilowanie wydajności}
Ocena wydajności będzie obejmować pomiar czasu inferencji w scenariuszach
jednordzeniowych i wielordzeniowych. Ważne jest nie tylko średnie FPS, ale
także analiza percentyli opóźnienia, które oddają stabilność działania.
Profilowanie zostanie przeprowadzone na sprzęcie referencyjnym (CPU) oraz
uzupełnione o pomiary FLOPs i zużycia pamięci. Takie podejście pozwala
jednoznacznie wskazać, czy poprawa jakości jest uzasadniona wzrostem kosztu.

\section{Funkcje straty i optymalizacja}
W segmentacji semantycznej standardowo stosuje się funkcje straty oparte o
entropię krzyżową, często z wagami klas w celu kompensacji nierównowagi
pomiędzy kategoriami. W projekcie rozważa się dodatkowo użycie strat pomocniczych,
np. strat na niższych rozdzielczościach, aby stabilizować uczenie wczesnych
warstw. Optymalizacja zostanie przeprowadzona przy użyciu standardowych
algorytmów (np. Adam lub SGD z momentum) z harmonogramem zmiany kroku uczenia.

\section{Ograniczenia i ryzyka}
Największym ryzykiem jest brak wystarczającej poprawy jakości w stosunku do
modeli bazowych przy zachowaniu niskiego kosztu obliczeniowego. Istnieje również
ryzyko, że mechanizm uwagi mimo uproszczeń nadal będzie generował znaczny narzut
czasowy. Dodatkowo, wyniki uzyskane na Cityscapes mogą nie w pełni uogólnić się
na inne domeny, co ogranicza przenaszalność rozwiązania.

\section{Struktura docelowej pracy}
Planowana praca końcowa będzie składała się z rozdziału wprowadzającego,
przeglądu literatury, opisu architektury FAscnn++, części eksperymentalnej oraz
wniosków. Sprawozdanie stanowi podstawę do rozwinięcia pierwszych rozdziałów
i stworzenia spójnej narracji, która zostanie uzupełniona o wyniki eksperymentalne.

\section{Harmonogram dalszych prac}
Na dalszym etapie planuje się przeprowadzenie pełnych eksperymentów na
Cityscapes, w tym treningu wariantów FAscnn++ oraz modeli bazowych. Kolejny krok
to analiza wyników, przygotowanie tabel i wykresów oraz opracowanie wniosków
z eksperymentów ablacjnych. Ostatni etap obejmie dopracowanie opisu architektury,
integrację wyników z częścią teoretyczną oraz przygotowanie pełnej wersji pracy.

\chapter{Przegląd literatury}
Przegląd literatury obejmuje cztery obszary: (1) klasyczne modele segmentacji
oparte na CNN, (2) lekkie i szybkie architektury real-time, (3) mechanizmy uwagi
i transformatory wizualne oraz (4) zbiory danych i standardy ewaluacji.

\section{Klasyczne modele segmentacji}
Podstawowym kamieniem milowym są sieci FCN, które jako pierwsze pokazały, że
konwolucyjne modele można adaptować do predykcji per-pikselowej
\parencite{long2015fully}. Kolejnym krokiem była architektura U-Net, bardzo
popularna w segmentacji medycznej, wykorzystująca połączenia typu skip
\parencite{ronneberger2015u}. SegNet zaproponował podejście z enkoderem i
dekoderem opartym o indeksery maksymalnego poolingu
\parencite{badrinarayanan2017segnet}. Z kolei DeepLab zastosował konwolucje
atrous i CRF w celu poprawy detali granic \parencite{chen2018deeplab}.
PSPNet wprowadził pyramid pooling dla lepszego modelowania kontekstu globalnego
\parencite{zhao2017pspnet}, a ICNet wykorzystał wieloskalowe przetwarzanie dla
przyspieszenia segmentacji \parencite{zhao2018icnet}.

\section{Wieloskalowość i fuzja cech}
Wiele architektur opiera się na fuzji informacji z różnych poziomów
rozdzielczości. Feature Pyramid Networks (FPN) wprowadziły ujednolicony sposób
łączenia map cech w detekcji, ale podejście to znalazło również zastosowanie w
segmentacji \parencite{lin2017fpn}. W praktyce fuzja wieloskalowa pozwala
zachować drobne detale z płytszych warstw i łączyć je z semantyką z warstw
głębokich. Podejścia takie stanowią inspiracje dla projektowania multi-branch
w FAscnn++, gdzie każda gałąź odpowiada za inny zakres skali.

\section{Lekkie architektury real-time}
ENet był jedną z pierwszych architektur zaprojektowanych dla segmentacji w
czasie rzeczywistym, z naciskiem na małe zużycie pamięci i parametrów
\parencite{paszke2016enet}. Fast-SCNN rozszerzył te idee o efektywne moduły
downsamplingu i fuzji cech \parencite{mehta2019fast}. BiSeNet zaproponował dwie
gałęzie przetwarzania: jedna szybka i druga dokładna, łączone w module fuzji
\parencite{yu2018bisenet}. W kontekście lekkich sieci istotne są również prace
dotyczące głębokościowo-separowalnych konwolucji (MobileNet, Xception), które
redukują koszty obliczeniowe \parencite{howard2017mobilenets,chollet2017xception}.

Typowe strategie w modelach real-time obejmują wczesne zmniejszanie rozdzielczości,
wykorzystanie bloków typu bottleneck oraz ograniczenie liczby kanałów w głębszych
warstwach. W praktyce prowadzi to do znacznego spadku FLOPs kosztem pewnej utraty
detali. Kluczowe jest więc dobranie takich mechanizmów fuzji, które pozwolą
odzyskiwać szczegóły w dekoderze bez znacznego narzutu obliczeniowego.

\section{Uwaga i transformatory}
Mechanizmy uwagi zrewolucjonizowały przetwarzanie sekwencji i znalazły zastosowanie
w wizji komputerowej \parencite{vaswani2017attention,dosovitskiy2020vit}.
Transformatory wizualne, takie jak SegFormer, potrafią modelować zależności
globalne bez kosztownych dekoderów \parencite{xie2021segformer}.
W kontekście lekkich modeli ważne są prace nad uwagą liniową, która redukuje
złożoność obliczeniową, np. Linformer oraz Performer
\parencite{wang2020linformer,choromanski2020performer}.

Rozwijane są również podejścia hybrydowe, które łączą konwolucje z uwagą.
Pozwalają one wykorzystać lokalną indukcję konwolucyjną, zachowując jednocześnie
zdolność do modelowania relacji dalekiego zasięgu. W praktyce takie podejścia
mogą być szczególnie atrakcyjne dla architektur real-time, ponieważ ograniczają
koszt uwagi do wybranych warstw lub skompresowanych reprezentacji.

\section{Zbiory danych i metryki}
Cityscapes to jeden z najważniejszych zbiorów danych do segmentacji miejskiej,
obejmujący wysokiej jakości adnotacje pikselowe
\parencite{cordts2016cityscapes}. Standardowymi metrykami oceny są mIoU oraz
Pixel Accuracy, uzupełniane o miary wydajności (FPS, opóźnienie).
W literaturze spotyka się również porównania z innymi zbiorami, takimi jak
CamVid, ADE20K czy COCO-Stuff \parencite{brostow2008camvid,zhou2017ade20k,caesar2018cocostuff},
jednak Cityscapes pozostaje dominującym wyborem w przypadku architektur real-time.

Zbiory miejskie, takie jak Cityscapes i CamVid, charakteryzują się dużą liczbą
klas związanych z infrastrukturą drogową, co sprawia, że modele muszą dobrze
rozróżniać obiekty o podobnym wyglądzie, np. elewacje, pojazdy i chodniki.
Z kolei zbiory ogólne, takie jak ADE20K czy COCO-Stuff, obejmują większą
różnorodność scen, co pozwala testować uogólnianie modeli poza domeną ruchu
drogowego. W niniejszej pracy Cityscapes stanowi główny punkt odniesienia,
ponieważ odpowiada celom segmentacji w scenach miejskich.

\section{Segmentacja instancyjna i panoptyczna}
Choć zakres pracy skupia się na segmentacji semantycznej, warto wskazać prace
związane z segmentacją instancyjną i panoptyczną. Mask R-CNN zaproponował
uniwersalny framework do jednoczesnej detekcji i segmentacji instancji
\parencite{he2017maskrcnn}. Koncepcja segmentacji panoptycznej łączy segmentację
semantyczną i instancyjną w jednym ujednoliconym zadaniu
\parencite{kirillov2019panoptic}. Te kierunki pokazują, jak ważna staje się
kompletna interpretacja sceny, ale jednocześnie podkreślają rosnącą złożoność
modeli i potrzebę optymalizacji wydajności.

\section{Kompresja i optymalizacja modeli}
W praktyce istotnym elementem jest kompresja modeli. Techniki takie jak
przycinanie wag (pruning), kwantyzacja i destylacja wiedzy pozwalają znacznie
zredukować rozmiar modelu i przyspieszyć inferencję
\parencite{han2015deepcompression,jacob2018quantization,hinton2015distilling}.
Choć celem projektu FAscnn++ jest zaprojektowanie lekkiej architektury, wykorzystanie
tych technik może dodatkowo poprawić wydajność w docelowym wdrożeniu.

Warto podkreślić, że kompresja modeli jest komplementarna wobec projektowania
architektury. Nawet najlepsze techniki kompresji nie zrekompensują nadmiernej
złożoności bazowej sieci, dlatego w projekcie najpierw zakłada się optymalizację
struktury, a dopiero w kolejnym kroku rozważa się kwantyzację lub pruning.
Taki kierunek pozwala zachować lepszą stabilność trenowania i kontrolować wpływ
poszczególnych zabiegów na jakość segmentacji.

\chapter{Podsumowanie postępów}
Dotychczas wykonano analizę literatury, wybrano zbiór Cityscapes oraz
zaimplementowano pipeline treningowy i testowy. Dodatkowo wdrożone zostały
modele bazowe (ENet i Fast-SCNN) oraz kilka wariantów architektury FAscnn++.
Kolejnym krokiem jest dopracowanie mechanizmu uwagi i przeprowadzenie
szczegółowej ewaluacji porównawczej wraz z eksperymentami ablacjnymi.

\printbibliography

\end{document}
